\section{What we compare ourselves to, and where it enters the float}

Requirement R3 is the one most easily faked. A related-work paragraph naming eight papers is not a
comparison; a table in which a named prior procedure occupies a row and gets a number is. The list
below is every piece of prior work that can occupy a row or a column of a Table~1 on this object,
what it becomes when it does, and whether we have run it.

\begin{center}\footnotesize
\begin{tabular}{@{}p{3.3cm}p{3.6cm}p{3.9cm}p{5.2cm}@{}}
\toprule
prior work & what it is & where it enters a Table~1 & measured? \\
\midrule
Perkovi\'c et al., generalised adjustment criterion; Henckel et al., optimal set &
the construction that turns a graph into the committed set &
the estimator, identical in every row of every candidate; also the audit oracle &
\yes\ and it is a shared defect to fix: both sides' validity checks must use this criterion
rather than the back-door one, and the two disagree on a handful of committed sets \\
\addlinespace
Guo and Perkovi\'c, enumeration in an MPDAG (arXiv:2010.08611) &
the interval over every adjustment set the class admits &
the row that discards the knowledge, and the denominator of the calibrated width &
\yes\ every cell \\
\addlinespace
Taeb, Guo and Henckel, model-oriented graph distances (arXiv:2511.10625) &
the order and the distance the whole object is built on &
the instrument that counts elementary graph revisions, one rung below ours &
\yes\ 3{,}014 rows \\
\addlinespace
structural Hamming distance &
the standard graph-distance baseline &
a column of the ranking table &
\yes\ synthetic arm only \\
\addlinespace
number of asserted claims $|\Kset|$ &
the naive instrument an analyst already has &
a rung of the certificate ladder and a column of the ranking table &
\yes\ both arms \\
\addlinespace
treatment-to-set separation &
a free structural rule read off the analyst's graph &
the cheapest rung of the certificate ladder &
\yes\ 3{,}014 rows \\
\addlinespace
leave-one-out and leave-two-out re-closure screen &
not in the literature, and the first thing a competent reviewer proposes &
a row of the price table and of the policy table &
\yes\ in the policy panel; \no\ missing from the ranking table, and owed \\
\addlinespace
Cinelli and Hazlett, omitted-variable bound &
the sensitivity analysis the field actually runs &
one row at its own oracle parameter, with the count of rows where it is defined &
\pt\ defined on 666 of 808 rows; covers 0.964 at 0.445 of the class width against 0.632 for
ours on the same rows; on the 101 rows where the committed set contains a descendant of the
treatment it covers 0.134 \\
\addlinespace
Rosenbaum; Masten and Poirier, breakdown frontiers &
the conceptual analogue and the reporting form &
the shape of the report table, not a computable baseline on this object &
\no\ by construction \\
\bottomrule
\end{tabular}
\end{center}

Two things follow. The first is that a candidate satisfies R3 only if a named prior procedure gets
a number in the same float, not a mention in the caption; by that test the price table and the
certificate ladder pass, the report table and the funnel do not. The second is that the Cinelli and
Hazlett row is the only comparison against a sensitivity analysis the field runs, and it is
favourable to them, because at their own oracle parameter their bound is narrower than our interval.
Printing it is what makes the table an argument rather than an advertisement, and the honest reading
is that it covers by assuming a model that a reversed orientation sits outside: on the rows where
the committed set contains a descendant of the treatment, the same formula at the same oracle covers
about one row in seven.
